\scanpage{249}
\mspage{123}
% archive sheet number: 248

$\underline{1}^{\circ}$) Si $f:X\unsure{-|\to}Y$ est dans $W$, alors il l'est couniversellement \unsure{sous|sur} $\varnothing$, i.e. pour \unsure{tt|$\forall$} $Z$ dans $\mathrm{Cat}$
\[
 f\amalg\mathrm{id}_Z:X\amalg Z\longrightarrow Y\amalg Z\quad\text{est dans }W.
\]

(\uline{NB} Ceci \deleted{\unsure{signifie que}}\inserted{équivaut à :} \unsure{tte somme finie|ttes sommes finies} de flèches $f_i\in W$ est dans $W$ --- et c'est $=$ pour sommes infinies, si Loc $(7)$, axiome des limites, est satisfait.)

$2^\circ$) Soit $Y_0\hookrightarrow Y$ une inclusion d'un \unsure{summand|sommand} direct d'un objet de $\mathrm{Cat}$. Alors si $X\xrightarrow{f}Y$ est dans $W$, de $=$ $X_0=X\mid Y_0-Y_0$.

(\uline{NB} Ceci équivaut à la condition que l'inclusion d'u\deleted{\unsure{ne}}\unsure{n} summand direct \deleted{dans $Y$} est toujours une $W$-fibration)

De plus, si $\pi_0(Y)$ est fini, ou si $W$ satisfait Loc $(7)$\inserted{(l'axiom. des limites)}, il suffit \deleted{dans $2^\circ$ \unsure{que}} d'exiger $2^\circ$) pour $Y_0$ une composante connexe de $Y$.

\inserted{c)} \deleted{\uline{Corollaire} 9.5. Soit $(f_i)_{i\in I}$ une} \inserted{Pour tte} famille \uline{finie} \inserted{$(f_i)_{i\in I}$} de flèches \deleted{de $\mathrm{Cat}$. Quand $W$ satisfait} $f_i:X_i-Y_i$ de $\mathrm{Cat}$, \deleted{Pour que} on a
\[
 \coprod f_i\in W\Longrightarrow f_i\in W\quad\text{pour tt }i\in I.
\]

De plus, si $W$ satisfait Loc $(7)$, alors il est inutile dans c) de se borner à $I$ fini.

(Et en tous cas, les conditions équivalentes a) b) c) impliquent pour \uline{tte} famille de flèches $(f_i)_{i\in I})$\unsure{;|,} \deleted{\unsure{que}} $\coprod f_i\in W\Rightarrow f_i\in W$ $\forall i\in I$.)
